\section{Local force-escorted work distillation}\label{app:local-work}
The physical teacher combines local force-guided proposals with energy-weighted
flow matching \citep{dern2025energyweighted,fab}. Normalized anchor references
make the escorted-work weights explicit \citep{escorted2008}.

\paragraph{Energy, force and mirror symmetry.}
The energy model assigns a scalar $E(X)$ to a geometry and provides its force
$F(X)=-\nabla_X E(X)$, a vector for each atom. The force points in the direction
of decreasing energy. We average a geometry and its inversion through the origin:
$E_+(X)=[E(X)+E(-X)]/2$ and $F_+(X)=[F(X)-F(-X)]/2$.
The opposite sign in the force average preserves $F_+=-\nabla E_+$.
This construction assigns equal energy to mirror-related geometries and opposite
forces to their inverted coordinates. We remove the mean atom force before
constructing displacements, keeping them in the centered coordinate space.

\paragraph{Local reference and target.}
For a generated training anchor $A$ with accepted perceived graph $G_A$, define
$q_A(X)=\mathcal N_{\HH_N}(A,\sigma_A^2I)$.
Write $kT=k_B T_{\rm phys}=0.0258519998$ eV for the thermal energy at 300 K.
We set
\begin{equation}
 \sigma_A=\min\left(0.03\,\text{\AA},
 \sqrt{\frac{0.10\,\text{\AA}\,kT}{\|F_+(A)\|}}\right),\qquad
 \delta_A=\frac{\sigma_A^2}{kT}F_+(A).
\end{equation}
A vanishing force selects the maximal width. The force norm is the Frobenius
norm over all atom-force components. These choices bound that same norm of the
displacement by $0.10$\,\AA\ and the Gaussian width by $0.03$\,\AA.
These quantities are fixed before
drawing $X\sim q_A$. The escort $Y=X+\delta_A$ is a translation on the centered
subspace and has determinant one. The local unnormalized target is
\begin{equation}
 \rho_A(Y)=q_A(Y)\exp\!\left[-\frac{E_+(Y)-E_+(A)}{kT}\right]
 \mathbf1_{G_A}(Y),
\end{equation}
where the support requires the same labelled perceived bond orders, formal charges
and radical counts, together with the connected/nonoverlap assay. The graph is
inferred from the anchor coordinates. The resulting target is a local energy tilt
restrained around $A$ by $q_A$.

\paragraph{Why the force displacement has this form.}
Locally approximate the molecular energy by its value and force at the anchor:
$E_+(Y)-E_+(A)\simeq-\langle F_+(A),Y-A\rangle_F$, where the inner product sums
force times displacement over every atom and coordinate. Completing the square
in the Gaussian gives
\begin{equation}
 q_A(Y)\exp\!\left[\frac{\langle F_+(A),Y-A\rangle_F}{kT}\right]
 \ \propto\
 \mathcal N_{\HH_N}\!\left(Y;A+\frac{\sigma_A^2}{kT}F_+(A),\sigma_A^2I\right).
 \label{eq:linear-boltzmann-escort}
\end{equation}
Thus the force escort samples the local Boltzmann tilt exactly for a linearized
energy, before applying the same-graph restriction. The Gaussian restraint acts
like a virtual harmonic spring with stiffness $kT/\sigma_A^2$: the force shifts
its mean by force divided by stiffness. This explains the displacement scale.
The graph restriction subsequently truncates this local distribution.

\paragraph{Complete local work.}
With reference Hamiltonian $H_0=-kT\log q_A$ and target Hamiltonian
$H_1=H_0+E_+-E_+(A)$ plus hard support, the generalized quench work is
\begin{equation}
 W_A(X)=E_+(Y)-E_+(A)+kT[\log q_A(X)-\log q_A(Y)]
\end{equation}
on valid support, and $+\infty$ otherwise. Change of variables gives
$\mathbb E_{q_A}[\exp(-W_A/kT)]=\int\rho_A(Y)\,dY$.
The two reference-density terms correct for the escort's proposal shift.
Analytic checks recover constant work for a constant-force escort and the
normalizer and weighted mean for a quadratic energy.
For the linearized energy in Equation~\eqref{eq:linear-boltzmann-escort},
substitution gives the same work
$W_A=-\sigma_A^2\|F_+(A)\|_F^2/(2kT)$ for every supported candidate.
Uniform weights are therefore appropriate in this approximation. Complete-work
weights use the actual nonlinear energy to correct deviations from it. This
distinguishes the inexpensive force-only teacher from its energy-weighted extension.

\paragraph{Training and controls.}
Each continuation selects eight distinct 17--28-atom training compositions by a
fixed metadata hash and generates 32 anchors per composition. These compositions
are disjoint from evaluation. Invalid anchors remain in the generation counts
and are excluded from the teacher pool. Each valid anchor produces eight escorted particles. Particles outside
its declared support receive zero weight. The work arm samples training endpoints
using self-normalized $\exp(-W_A/kT)$ weights; the unweighted escort arm samples the
same valid particles uniformly. A replay arm uses the retained raw anchors, and a
frozen model supplies the reference. All-zero-particle anchors are recorded and
excluded identically from all student pools.

Each student alternates 500 original-reference FM updates with 500 generated-anchor
endpoint updates, sharing initialization, anchor draws, optimization and sampling
budgets. Physical preparation costs differ: replay needs no force queries, the
unweighted escort needs anchor forces, and complete work additionally needs
energies at valid proposed endpoints. Actual shared research cost and these
method requirements are reported separately. The final generator uses the
learned network alone. Its temperature feature stays fixed throughout training
and evaluation; 300 K determines the local teacher weights. Self-normalization
introduces finite-particle bias, and the reported teacher ESS is computed over
the eight local particles at each anchor.
Effective sample size (ESS) is $(\sum_j w_j)^2/\sum_j w_j^2$ for candidate weights
$w_j$. It equals the candidate count for equal weights and approaches one when
a single candidate dominates.

\paragraph{Direct fine-tuning results.}
We train three students for 1,000 updates in each run. Across 512 generated
anchors, 351 pass the graph checks and 349 have at least one allowed force-shifted
candidate. Composition-level mean ESS is about 1.27--1.87 out of 8. The weights
are therefore concentrated on few candidates, limiting resolution of the local
target with this pool size.

\begin{table}[t]
\centering
\caption{Graph-valid outputs after direct physical fine-tuning, out of 640
attempts per run. The frozen model receives no updates; all other variants have
the same additional training budget.}
\begin{tabular}{lrrrr}
\toprule
Run & Frozen & Replay & Unweighted escort & Complete work\\
\midrule
1 &418&340&330&344\\
2 &377&345&371&362\\
\bottomrule
\end{tabular}\label{tab:local-work-direct}
\end{table}

Among paired outputs that are both graph-valid, complete-work training gives
lower energy than unweighted-escort training by 0.00758 eV/atom (paired 95\%
interval $[-0.01424,-0.00113]$ for the signed difference). Force RMS is lower by
0.16174 eV/\AA\ ($[-0.27660,-0.05590]$). Both runs favor complete work, and
the all-output comparison has the same direction.

However, all three directly fine-tuned students perform worse than the frozen
model. Complete work has 55.16\% graph validity versus 62.11\% for the frozen
model, and energy on paired, jointly valid outputs is \emph{higher} by
0.10587 eV/atom ($[0.09009,0.12168]$). Replay also worsens performance after
training on the mixed targets. These results motivate subtracting the matched
replay update rather than using the physical student directly.

Teacher preparation uses 6,232 raw energy/force queries: 702 at anchors and
5,530 at allowed candidates. The unweighted escort could omit candidate-energy
queries, and replay could omit all physical queries. Scoring the outputs of the
four models uses a further 10,240 queries, separate from generation cost.


\paragraph{Applying the paired physical update.}
Instead of deploying the directly fine-tuned student, we use parameter arithmetic
\citep{ilharco2023arithmetic} to subtract its matched replay update:
\begin{equation}
 \theta_{\rm final}=\theta_{\rm base}
   +(\theta_{\rm physical}-\theta_{\rm replay}),
\end{equation}
for both unweighted-escort and complete-work students. The coefficient is fixed
at one. We change parameter tensors only; buffers and the source distribution
remain fixed.

\begin{table}[t]
\centering
\caption{Paired updates evaluated after the direct fine-tuning experiment.
Energy changes use paired, jointly valid outputs and are relative to the frozen
model, with paired 95\% confidence intervals. Counts show the two generation
streams.}
\begin{tabular}{lrr}
\toprule
Model & Graph-valid outputs (two streams) & Energy change (eV/atom)\\
\midrule
Frozen &418/640; 377/640&$0$\\
Paired force update &407/640; 387/640&$-0.01554\;[-0.02084,-0.01031]$\\
Paired work update &406/640; 382/640&$-0.01604\;[-0.02188,-0.01013]$\\
\bottomrule
\end{tabular}\label{tab:thermal-transfer}
\end{table}

Both paired variants lower energy and force RMS in both runs. For complete work,
force RMS on paired, jointly valid outputs falls by 0.38403 eV/\AA\ (interval
$[-0.48124,-0.28932]$); all-attempt energy and force scores also improve. Its
graph-validity difference from the frozen model is $-0.55$ percentage points,
with interval $[-2.58,1.56]$.

The complete-work and unweighted paired variants have similar scores. Complete
work minus unweighted gives $-0.00243$ eV/atom on jointly valid pairs
($[-0.00756,0.00264]$) and $+0.00022$ eV/atom over all attempts
($[-0.00557,0.00649]$). Force and graph-validity intervals also include zero.
Teacher preparation requires 702 anchor queries for the unweighted variant
versus 6,232 for complete work. Both retain the 128-call neural sampler.

Evaluating these two paired variants adds 2,560 generated outputs and 5,120
physical queries. Combining parameters requires no optimizer steps. Based on
these development results, we select the force-only update for the
new-composition evaluation in Section~\ref{sec:fresh-physics}.

